\section{はじめに}
近年, ストリーミングサービスなどメディアコンテンツに関連するサービスが大きな広がりを見せており, 社会におけるメディアコンテンツの流動性が増している. 日々大量のコンテンツが消費され, 次から次へと新しいものが求められるようになった. 

このようなコンテンツ需要の高まりに応えるストーリー創作の困難は大きい. 多くのコンテンツは, ストーリーを骨格として形成され, その新奇性と多様性がコンテンツの魅力と供給に直結するためである. ストーリーの創作には高度な創造性を要し, コンテンツ需要が加速する中, 十分な新しさや多様性をもつ高品質なストーリーを, 社会の消費に見合う速度で生み出すことに大変な労力が必要になっているのである.

このような負担を軽減するため, 創造性を支援する技術が望まれており, AI技術の応用が期待されている.  

創造性とは, 価値のある新たな発想を生み出す能力のことであり, 価値には, 面白さや美しさ, 有益さといったさまざまな意味が含まれる\cite{ref:Boden09}. したがって, ストーリー創作における創造性の支援は, これまでにない面白さをもつ物語への気づきを与えるようなものであることが望まれる. 

従来のAIがこなせるタスクは単純なものに限られていたため, このような創造性を必要とするタスクに適用することは難しいという議論もあった\cite{ref:Boden09}. 

しかし現在, より高度で幅広いタスクにも対応可能な Foundation Model\cite{ref:Bommasani21}が登場し, AI の適用範囲は創造的活動にまで広がっていると注目されている\cite{ref:Frich19}. 特に, GPT-4\cite{ref:OpenAI23} をはじめとした大規模言語モデル(LLM) への期待は大きい. 

今日の LLM は, 自然言語処理におけるベンチマークにおいて驚異的な数値を叩き出す\cite{ref:Zhao23}. 自然言語を理解しているかのように応答し, 生成される文章は人が書いた文章に匹敵するほどである. 
%自然言語処理分野における様々なタスクで高い性能を発揮するようになっており\cite{ref:Zhao23}, ストーリー創作における創造性支援の可能性を大きく広げている.



様々な言語タスクに応用可能なLLMの有用さは, 創造性支援にも活用できると考えられ, これまでも, ストーリークリエーターの創造性支援に LLM を用いる試みが行われてきた\cite{ref:Mirowski23}.

しかしながら, その活用方法は未だ確立されていない. LLM から有用な出力を引き出すためにはプロンプトに工夫が必要であると知られており\cite{ref:Chen23}, AIの活用に馴染みのないストーリークリエーターでも, 負担なく有意義に LLM を活用することができる体系の構築が求められている. 

そこで我々は, 物語構造を活用し, プロのストーリークリエーターが LLM を用いたストーリー創作を効率的に行うことができる, ユーザーと LLM の仲介を担う御用聞き AI を web アプリとして開発した. このアプリケーションは画面の指示に従う直感的な操作で, 物語構造分析に基づいたプロット生成を実現する. さらに, 生成されたプロットの改善を目的とする, LLM との体系的でインタラクティブな共創手法を提供する. 我々はこのアプリケーションの開発により, LLM を活用したプロのストーリークリエーターの創作支援を探求した.

%そこで我々は, 物語構造を活用し, プロのストーリークリエーターが LLM を用いたストーリー創作を効率的に行うことができる web アプリケーションを開発した. ユーザーは, このアプリケーションの操作に複雑な入力を必要としない. 画面の指示に従い選択と入力を進めていくことで, 事前に分析された物語構造に沿うプロットの生成を行うことができる. さらに, このアプリケーションでは, 生成されたプロットの改善を目的とする, 体系的でインタラクティブな LLM とのやり取りを行うことができる. このアプリケーションの開発により, プロのストーリークリエーターの創作を支援する LLM の活用方法を模索した.

%研究の成果として、本研究チームは物語構造に基づくフレームワークを適用し、専門のクリエーターがLarge Language Models（LLM）を駆使してストーリーテリングを効率的に行うためのウェブベースのアプリケーションを開発した。このアプリケーションのユーザーインターフェースは複雑な操作を必要とせず、直感的に画面のガイダンスに従って選択肢を選び、入力を行うことで、事前に解析済みの物語構造を反映したプロット生成を容易に実現する。加えて、アプリケーションは生成されたプロットの洗練を促進する目的で、体系的かつ対話的なLLMとのインタラクションを可能にしている。このアプリケーションの導入により、プロフェッショナルなストーリークリエーターがLLMを活用する新たな方法論についての検討を行い、創造過程における支援を目指すものである。